\documentclass[aspectratio=169]{beamer}

% ---------- ENCODING & FONTS ----------
\usepackage[utf8]{inputenc}
\usepackage[T1]{fontenc}

% ---------- THEME & LOOK ----------
\usetheme{Madrid}
\usecolortheme{seahorse}
\usefonttheme{professionalfonts}
\setbeamertemplate{navigation symbols}{}
\setbeamercovered{transparent}
\setbeamertemplate{itemize items}[circle]
\setbeamertemplate{enumerate items}[default]
\setlength{\parskip}{0.25em}

% ---------- COLORS ----------
\definecolor{accent}{RGB}{0,92,184}
\definecolor{soft}{RGB}{233,244,255}
\definecolor{canadared}{RGB}{196,30,58}
\definecolor{leafgreen}{RGB}{0,130,80}
\definecolor{goldenrod}{RGB}{218,165,32}
\definecolor{deeporange}{RGB}{220,90,20}
\definecolor{shiftpurple}{RGB}{130,50,180}
\definecolor{tealblue}{RGB}{0,128,128}
\definecolor{lightred}{RGB}{245,210,210}
\setbeamercolor{structure}{fg=accent}
\setbeamercolor{title}{fg=white,bg=accent}
\setbeamercolor{frametitle}{fg=accent}
\setbeamercolor{block title}{bg=accent!12,fg=accent}
\setbeamercolor{block body}{bg=soft}

% ---------- PACKAGES ----------
\usepackage{booktabs}
\usepackage{amsmath, amssymb}
\usepackage{array}
\usepackage{multirow}
\usepackage{hyperref}
\usepackage{tikz}
\usepackage{pgfplots}
\pgfplotsset{compat=1.18}
\usetikzlibrary{arrows.meta, positioning, calc, shapes.geometric}
\hypersetup{colorlinks=true,linkcolor=accent,urlcolor=accent}

% ---------- SAFE TRANSITION FALLBACK ----------
\makeatletter
\@ifundefined{transfade}{\newcommand{\transfade}[1][]{}}{}
\makeatother

% ---------- TITLE ----------
\title{\textbf{Fiscal Policy}}
\subtitle{{\large ECON 101 -- Principles of Macroeconomics}}
\author{Muhebullah Karimzada}
\date{Spring 2026}

% ============================================================
\begin{document}
% ============================================================

\begin{frame}[plain]
  \titlepage
\end{frame}

\begin{frame}{Chapter Outline}
\transfade
\begin{itemize}[<+->]
  \item \textbf{12.1} What Is Fiscal Policy?
  \item \textbf{12.2} The Effects of Fiscal Policy on Real GDP and the Price Level
  \item \textbf{12.4} Government Purchases and Tax Multipliers
  \item \textbf{12.5} The Limits of Fiscal Policy as a Stimulus
  \item \textbf{12.6} Deficits, Surpluses, and Federal Government Debt
  \item \textbf{12.7} The Effects of Fiscal Policy in the Long Run
  \item \textbf{Appendix} A Closer Look at the Multiplier
\end{itemize}
\end{frame}

% ============================================================
\section{12.1 What Is Fiscal Policy?}
% ============================================================

\begin{frame}{12.1 Learning Objective}
\transfade
\begin{block}{Goal}
Define fiscal policy and distinguish between automatic stabilizers and discretionary fiscal policy.
\end{block}
\begin{itemize}[<+->]
  \item \textbf{Fiscal policy:} changes in \alert{federal government taxes and spending} intended to achieve macroeconomic objectives.
  \begin{itemize}[<+->]
    \item Provincial spending is important but primarily aimed at service delivery, not national macrostabilization.
  \end{itemize}
  \item The federal government can \textbf{shift aggregate demand} by changing spending and tax decisions.
  \item \textbf{Canada 2020 example:} the federal government launched \$314 billion in emergency spending (CERB, CEWS, CERS) --- the largest fiscal intervention in Canadian history.
\end{itemize}
\end{frame}

\begin{frame}{Automatic Stabilizers vs. Discretionary Fiscal Policy}
\transfade
\begin{columns}
\column{0.48\textwidth}
\begin{block}{\textcolor{leafgreen}{Automatic Stabilizers}}
Spending and taxes that automatically change with the business cycle \textbf{without new legislation}.
\begin{itemize}
  \item Employment Insurance (EI): rises in recessions as more workers qualify.
  \item Income tax: falls in recessions (lower incomes $\Rightarrow$ lower tax revenue).
  \item Social assistance: rises when unemployment increases.
\end{itemize}
\textit{Act fast, no legislative delay.}
\end{block}
\column{0.48\textwidth}
\begin{block}{\textcolor{deeporange}{Discretionary Fiscal Policy}}
Intentional government actions requiring new legislation.
\begin{itemize}
  \item 2009 Economic Action Plan: \$40B infrastructure stimulus.
  \item 2020 CERB: \$2{,}000/month to unemployed Canadians.
  \item 2024 GST holiday: temporary sales tax cut.
  \item 2025 income tax cut proposal.
\end{itemize}
\textit{Takes time; subject to political process.}
\end{block}
\end{columns}
\end{frame}

% ============================================================
\section{12.2 Effects of Fiscal Policy on GDP and the Price Level}
% ============================================================

\begin{frame}{12.2 Learning Objective}
\transfade
\begin{block}{Goal}
Explain how fiscal policy affects aggregate demand and how it can stabilize the economy.
\end{block}
\begin{itemize}[<+->]
  \item \textbf{Two channels:}
  \begin{itemize}[<+->]
    \item Changes in government purchases ($G$): affects AD \alert{directly}.
    \item Changes in taxes ($T$): affects disposable income $\Rightarrow$ consumption $\Rightarrow$ AD \alert{indirectly}.
  \end{itemize}
  \item \textbf{Expansionary fiscal policy:} $\uparrow G$ or $\downarrow T$ $\Rightarrow$ AD shifts right.
  \item \textbf{Contractionary fiscal policy:} $\downarrow G$ or $\uparrow T$ $\Rightarrow$ AD shifts left.
\end{itemize}
\end{frame}

\begin{frame}{Figure 12.1 -- Expansionary Fiscal Policy (Closing a Recession)}
\begin{center}
\begin{tikzpicture}
\begin{axis}[
  width=0.80\textwidth, height=0.70\textheight,
  xlabel={Real GDP (\$billions)},
  ylabel={Price Level ($P$)},
  xmin=0, xmax=3000, ymin=80, ymax=200,
  xtick={500,1000,1500,2000,2500},
  ytick={90,110,130,150,170,190},
  grid=major, grid style={dotted,gray!40},
  tick label style={font=\small},
  label style={font=\small},
  title style={font=\small\bfseries},
  title={Expansionary Fiscal Policy: $\uparrow G$ or $\downarrow T$ Closes Recessionary Gap}
]
\addplot[shiftpurple, very thick] coordinates {(2000,80)(2000,200)};
\node[right, font=\tiny, color=shiftpurple] at (axis cs:2050,185) {LRAS};
\addplot[goldenrod, very thick, domain=200:2800] {0.05*x + 30};
\node[right, font=\tiny, color=goldenrod] at (axis cs:2650,163) {SRAS};
% AD1 (below potential)
\addplot[accent!60, very thick, dashed, domain=200:2800] {215 - 0.055*x}
  node[right, font=\tiny, color=accent!60] {$AD_1$ (recession)};
% AD2 (after stimulus)
\addplot[leafgreen, very thick, domain=200:2800] {240 - 0.055*x}
  node[right, font=\tiny, color=leafgreen] {$AD_2$ (stimulus)};
% Eq1: 215-0.055x=0.05x+30 => 185=0.105x => x=1762; P=118
\addplot[accent, only marks, mark=*, mark size=3.5pt] coordinates {(1762,118)};
\draw[dashed, accent] (axis cs:1762,118)--(axis cs:1762,0) node[below,font=\tiny,color=accent]{$Y_1$};
\draw[dashed, accent] (axis cs:1762,118)--(axis cs:0,118) node[left,font=\tiny,color=accent]{$P_1$};
% Eq2: 240-0.055x=0.05x+30 => 210=0.105x => x=2000; P=130
\addplot[leafgreen, only marks, mark=*, mark size=3.5pt] coordinates {(2000,130)};
\draw[dashed, leafgreen] (axis cs:2000,130)--(axis cs:2000,0) node[below,font=\tiny,color=leafgreen]{$Y^*=Y_2$};
\draw[dashed, leafgreen] (axis cs:2000,130)--(axis cs:0,130) node[left,font=\tiny,color=leafgreen]{$P_2$};
\draw[->, very thick, leafgreen] (axis cs:1300,165) -- (axis cs:1600,165)
  node[above, midway, font=\tiny, color=leafgreen] {$\uparrow G$ or $\downarrow T$};
\node[font=\tiny, fill=leafgreen!10, draw=leafgreen, rounded corners, align=center]
  at (axis cs:600,175) {\textbf{Canada 2020:}\\CERB + CEWS\\closed recessionary gap};
\end{axis}
\end{tikzpicture}
\end{center}
\end{frame}

\begin{frame}{Figure 12.2 -- Contractionary Fiscal Policy (Fighting Inflation)}
\begin{center}
\begin{tikzpicture}
\begin{axis}[
  width=0.80\textwidth, height=0.70\textheight,
  xlabel={Real GDP (\$billions)},
  ylabel={Price Level ($P$)},
  xmin=0, xmax=3000, ymin=80, ymax=200,
  xtick={500,1000,1500,2000,2500},
  ytick={90,110,130,150,170,190},
  grid=major, grid style={dotted,gray!40},
  tick label style={font=\small},
  label style={font=\small},
  title style={font=\small\bfseries},
  title={Contractionary Fiscal Policy: $\downarrow G$ or $\uparrow T$ Reduces Inflationary Gap}
]
\addplot[shiftpurple, very thick] coordinates {(2000,80)(2000,200)};
\node[right, font=\tiny, color=shiftpurple] at (axis cs:2050,185) {LRAS};
\addplot[goldenrod, very thick, domain=200:2800] {0.05*x + 30};
\node[right, font=\tiny, color=goldenrod] at (axis cs:2650,163) {SRAS};
% AD1 (above potential - inflationary)
\addplot[canadared!60, very thick, dashed, domain=200:2800] {260 - 0.055*x}
  node[right, font=\tiny, color=canadared!60] {$AD_1$ (inflationary)};
% AD2 (after austerity)
\addplot[accent, very thick, domain=200:2800] {240 - 0.055*x}
  node[right, font=\tiny, color=accent] {$AD_2$ (austerity)};
% Eq1: 260-0.055x=0.05x+30 => x=2190; P=139.5
\addplot[canadared, only marks, mark=*, mark size=3.5pt] coordinates {(2190,140)};
\draw[dashed, canadared] (axis cs:2190,140)--(axis cs:2190,0) node[below,font=\tiny,color=canadared]{$Y_1 > Y^*$};
\draw[dashed, canadared] (axis cs:2190,140)--(axis cs:0,140) node[left,font=\tiny,color=canadared]{$P_1$};
% Eq2: x=2000; P=130
\addplot[accent, only marks, mark=*, mark size=3.5pt] coordinates {(2000,130)};
\draw[dashed, accent] (axis cs:2000,130)--(axis cs:2000,0) node[below,font=\tiny,color=accent]{$Y^*$};
\draw[dashed, accent] (axis cs:2000,130)--(axis cs:0,130) node[left,font=\tiny,color=accent]{$P_2$};
\draw[->, very thick, canadared] (axis cs:1700,165) -- (axis cs:1400,165)
  node[above, midway, font=\tiny, color=canadared] {$\downarrow G$ or $\uparrow T$};
\node[font=\tiny, fill=lightred, draw=canadared, rounded corners, align=center]
  at (axis cs:700,175) {\textbf{Canada 1995:}\\Martin's austerity\\closed inflationary\\pressures};
\end{axis}
\end{tikzpicture}
\end{center}
\end{frame}

\begin{frame}{Table 12.1 -- Countercyclical Fiscal Policy in Canada}
\transfade
\begin{center}
\begin{tabular}{p{0.18\linewidth} p{0.18\linewidth} p{0.32\linewidth} p{0.25\linewidth}}
\toprule
\textbf{Problem} & \textbf{Policy Type} & \textbf{Government Action} & \textbf{Canadian Example} \\
\midrule
Recession & Expansionary & $\uparrow G$ or $\downarrow T$ $\Rightarrow$ AD shifts right & CERB/CEWS 2020 (\$314B) \\
\addlinespace
Rising inflation & Contractionary & $\downarrow G$ or $\uparrow T$ $\Rightarrow$ AD shifts left & 1995 Martin budget cuts \\
\bottomrule
\end{tabular}
\end{center}
\vspace{0.5em}
\begin{itemize}[<+->]
  \item \textbf{Note:} in Table 12.1 the original text references the U.S. Congress. The Canadian equivalent is Parliament --- specifically the House of Commons, which must pass all spending bills and tax changes.
  \item Countercyclical policy aims to \textbf{dampen} the business cycle, not eliminate it.
\end{itemize}
\end{frame}

% ============================================================
\section{12.4 Government Purchases and Tax Multipliers}
% ============================================================

\begin{frame}{12.4 Learning Objective}
\transfade
\begin{block}{Goal}
Explain how the government purchases and tax multipliers work.
\end{block}
\begin{itemize}[<+->]
  \item A \$1 increase in government spending raises equilibrium GDP by \alert{more than \$1}.
  \item Why?
  \begin{itemize}[<+->]
    \item Initial spending $\Rightarrow$ incomes rise $\Rightarrow$ consumption rises $\Rightarrow$ incomes rise further $\Rightarrow$ \ldots
    \item This chain is the \textbf{multiplier effect}.
  \end{itemize}
  \item \textbf{Government purchases multiplier:}
  \[
    \frac{\Delta Y}{\Delta G} = \frac{1}{1 - \text{MPC}}.
  \]
  \item \textbf{Tax multiplier:}
  \[
    \frac{\Delta Y}{\Delta T} = -\frac{\text{MPC}}{1 - \text{MPC}}.
  \]
\end{itemize}
\end{frame}

\begin{frame}{Figure 12.3 -- The Multiplier Effect in AD-AS}
\begin{center}
\begin{tikzpicture}
\begin{axis}[
  width=0.82\textwidth, height=0.68\textheight,
  xlabel={Real GDP (\$billions)},
  ylabel={Price Level ($P$)},
  xmin=0, xmax=3000, ymin=80, ymax=200,
  xtick={500,1000,1500,2000,2500},
  ytick={90,110,130,150,170,190},
  grid=major, grid style={dotted,gray!40},
  tick label style={font=\small},
  label style={font=\small},
  title style={font=\small\bfseries},
  title={Multiplier: \$10B Increase in $G$ Raises Equilibrium GDP by \$50B (MPC $= 0.8$)}
]
\addplot[shiftpurple, very thick] coordinates {(2000,80)(2000,200)};
\node[right, font=\tiny, color=shiftpurple] at (axis cs:2050,185) {LRAS};
\addplot[goldenrod, very thick, domain=200:2800] {0.05*x + 30};
\node[right, font=\tiny, color=goldenrod] at (axis cs:2650,163) {SRAS};
% AD1
\addplot[accent, very thick, dashed, domain=200:2800] {215 - 0.055*x}
  node[right, font=\tiny, color=accent] {$AD_1$};
% AD2 (shifted right by full multiplier effect = $50B)
% If $G rises by 10, multiplier=5 => $Y rises by 50
% AD1 at x=1762, P=118. AD2 at x=1762+50=1812, but with upward-sloping SRAS
% price also rises. Let's just show a clean shift.
\addplot[leafgreen, very thick, domain=200:2800] {232 - 0.055*x}
  node[right, font=\tiny, color=leafgreen] {$AD_2$ (after multiplier)};
% Eq1: x=1762; P=118
\addplot[accent, only marks, mark=*, mark size=3.5pt] coordinates {(1762,118)};
\draw[dashed, accent] (axis cs:1762,118)--(axis cs:1762,0) node[below,font=\tiny,color=accent]{$Y_1$};
% Eq2: 232-0.055x=0.05x+30 => 202=0.105x => x=1924; P=0.05*1924+30=126
\addplot[leafgreen, only marks, mark=*, mark size=3.5pt] coordinates {(1924,126)};
\draw[dashed, leafgreen] (axis cs:1924,126)--(axis cs:1924,0) node[below,font=\tiny,color=leafgreen]{$Y_2$};
\draw[<->, very thick, deeporange] (axis cs:1762,88) -- (axis cs:1924,88)
  node[above, midway, font=\tiny, color=deeporange] {$\Delta Y = \$162B$};
\node[font=\tiny, fill=soft, draw=deeporange, rounded corners, align=center]
  at (axis cs:800,175) {$\Delta G = \$10B$ autonomous\\Multiplier $= 5$ (MPC $= 0.8$)\\BUT SRAS slopes up:\\part shows as $\uparrow P$};
\end{axis}
\end{tikzpicture}
\end{center}
\end{frame}

\begin{frame}{Government Purchases Multiplier vs. Tax Multiplier}
\transfade
\begin{columns}
\column{0.48\textwidth}
\begin{block}{\textcolor{accent}{Government Purchases Multiplier}}
\[
  \frac{\Delta Y}{\Delta G} = \frac{1}{1-\text{MPC}}
\]
\begin{itemize}
  \item \$1 of $G$ \textbf{directly} adds \$1 to AD.
  \item Then the induced consumption chain kicks in.
  \item With MPC $= 0.8$: multiplier $= 5$.
  \item $\Delta G = \$10B \Rightarrow \Delta Y = \$50B$.
\end{itemize}
\end{block}
\column{0.48\textwidth}
\begin{block}{\textcolor{canadared}{Tax Multiplier}}
\[
  \frac{\Delta Y}{\Delta T} = -\frac{\text{MPC}}{1-\text{MPC}}
\]
\begin{itemize}
  \item Tax cut raises disposable income.
  \item Households \textbf{save MPS} portion, spend MPC portion.
  \item First-round effect: \$MPC (not \$1).
  \item With MPC $= 0.8$: tax multiplier $= -4$.
  \item $\Delta T = -\$10B \Rightarrow \Delta Y = +\$40B$.
\end{itemize}
\end{block}
\end{columns}
\vspace{0.4em}
\begin{itemize}[<+->]
  \item The \textbf{government purchases multiplier is larger} (absolute value) than the tax multiplier: spending injection is fully spent, while a tax cut is partly saved.
\end{itemize}
\end{frame}

% ============================================================
\section{12.5 The Limits of Fiscal Policy}
% ============================================================

\begin{frame}{12.5 Learning Objective}
\transfade
\begin{block}{Goal}
Discuss the difficulties that can arise in implementing fiscal policy.
\end{block}
\begin{itemize}[<+->]
  \item Fiscal policy is powerful in theory but faces practical limitations:
  \begin{enumerate}[<+->]
    \item \textbf{Timing problems} (legislative and implementation lags).
    \item \textbf{Crowding out} (higher government borrowing raises interest rates, reducing private spending).
    \item \textbf{Political constraints} (fiscal policy is a political process, not just an economic one).
  \end{enumerate}
\end{itemize}
\end{frame}

\begin{frame}{Timing Problems}
\transfade
\begin{enumerate}[<+->]
  \item \textbf{Recognition lag:}
  \begin{itemize}[<+->]
    \item GDP data arrives 2--3 months after the quarter ends. By the time a recession is recognized, it may already be deep.
  \end{itemize}
  \item \textbf{Legislative delay:}
  \begin{itemize}[<+->]
    \item Parliament must debate and approve spending or tax changes.
    \item \textbf{Canadian example (2009):} Conservative minority government's Economic Action Plan took months to implement after the recession had begun.
  \end{itemize}
  \item \textbf{Implementation delay:}
  \begin{itemize}[<+->]
    \item Infrastructure projects take months to years to actually begin (design, tender, construction).
    \item \textbf{``Shovel-ready'' projects:} a political priority in 2009 precisely because of this lag.
  \end{itemize}
  \item Because of these delays, fiscal stimulus may arrive \alert{after the recession has ended} --- adding fuel to the next boom.
\end{enumerate}
\end{frame}

\begin{frame}{Crowding Out}
\transfade
\begin{block}{Crowding Out}
A decline in private expenditures (C, I, NX) that results from an increase in government purchases.
\end{block}
\begin{itemize}[<+->]
  \item \textbf{Mechanism:}
  \begin{itemize}[<+->]
    \item Government increases spending $\Rightarrow$ borrows more (issues more bonds).
    \item Higher bond supply $\Rightarrow$ bond prices fall $\Rightarrow$ interest rates rise.
    \item Higher interest rates $\Rightarrow$ lower private investment, lower consumer credit, stronger CAD (lower NX).
    \item Private spending falls --- partially offsetting the stimulus.
  \end{itemize}
  \item \textbf{Canadian context:}
  \begin{itemize}[<+->]
    \item Canada's federal debt is $\approx$\$1.3 trillion (2024).
    \item Higher debt-to-GDP ratios may gradually raise borrowing costs.
    \item But with BoC independence, crowding out is partially offset by monetary policy.
  \end{itemize}
\end{itemize}
\end{frame}

\begin{frame}{Figure 12.4 -- Crowding Out in the Loanable Funds Market}
\begin{center}
\begin{tikzpicture}
\begin{axis}[
  width=0.78\textwidth, height=0.70\textheight,
  xlabel={Quantity of Loanable Funds (\$billions)},
  ylabel={Real Interest Rate (\%)},
  xmin=0, xmax=10, ymin=0, ymax=8,
  xtick={0,2,4,6,8,10}, ytick={0,1,2,3,4,5,6,7,8},
  grid=major, grid style={dotted,gray!40},
  tick label style={font=\small},
  label style={font=\small},
  title style={font=\small\bfseries},
  title={Government Borrowing Raises Interest Rates, Crowding Out Private Investment}
]
% Demand for loanable funds
\addplot[canadared, very thick, domain=0:9] {-0.6*x + 7.7}
  node[right, font=\tiny, color=canadared] {\textbf{D} (private investment)};
% S1 (before deficit)
\addplot[leafgreen, very thick, domain=0:9] {0.6*x + 0.5}
  node[right, font=\tiny, color=leafgreen] {$S_1$ (private saving)};
% S2 (after deficit - supply reduced because govt is also borrowing)
\addplot[goldenrod, very thick, dashed, domain=0:8.5] {0.6*x + 2.2}
  node[right, font=\tiny, color=goldenrod] {$S_2$ (after govt deficit)};
% Eq1: -0.6x+7.7=0.6x+0.5 => 7.2=1.2x => x=6; P=4.1
\addplot[leafgreen, only marks, mark=*, mark size=3.5pt] coordinates {(6.0,4.1)};
\draw[dashed, leafgreen] (axis cs:6,4.1)--(axis cs:6,0) node[below,font=\tiny,color=leafgreen]{$Q_1$};
\draw[dashed, leafgreen] (axis cs:6,4.1)--(axis cs:0,4.1) node[left,font=\tiny,color=leafgreen]{$r_1$};
% Eq2: -0.6x+7.7=0.6x+2.2 => 5.5=1.2x => x=4.58; P=4.96
\addplot[goldenrod, only marks, mark=*, mark size=3.5pt] coordinates {(4.58,4.95)};
\draw[dashed, goldenrod] (axis cs:4.58,4.95)--(axis cs:4.58,0) node[below,font=\tiny,color=goldenrod]{$Q_2$};
\draw[dashed, goldenrod] (axis cs:4.58,4.95)--(axis cs:0,4.95) node[left,font=\tiny,color=goldenrod]{$r_2$};
\draw[<->, canadared, very thick] (axis cs:4.58,0.4) -- (axis cs:6.0,0.4)
  node[above, midway, font=\tiny, color=canadared] {Crowding out};
\node[font=\tiny, fill=lightred, draw=canadared, rounded corners, align=center]
  at (axis cs:2,6.8) {Govt deficit\\shifts S left\\$\Rightarrow r$ rises\\$\Rightarrow$ fewer private\\loans};
\end{axis}
\end{tikzpicture}
\end{center}
\end{frame}

% ============================================================
\section{12.6 Deficits, Surpluses, and Federal Government Debt}
% ============================================================

\begin{frame}{12.6 Learning Objective}
\transfade
\begin{block}{Goal}
Define budget deficit and federal government debt, and explain how the budget serves as an automatic stabilizer.
\end{block}
\begin{itemize}[<+->]
  \item \textbf{Budget deficit:} expenditures $>$ tax revenues.
  \item \textbf{Budget surplus:} expenditures $<$ tax revenues.
  \item \textbf{Federal government debt (national debt):} the total stock of outstanding Government of Canada bonds.
  \item \textbf{Canadian fiscal history:}
  \begin{itemize}[<+->]
    \item Deficits almost every year since 1970.
    \item One stretch of surpluses: 1998--2008 (Chrétien/Martin era).
    \item \alert{2020--21}: deficit of \$314 billion --- by far the largest ever.
    \item \alert{2024--25}: deficit $\approx$ \$40 billion.
  \end{itemize}
\end{itemize}
\end{frame}

\begin{frame}{Figure 12.5 -- Canada's Federal Budget Balance (1990--2024)}
\begin{center}
\begin{tikzpicture}
\begin{axis}[
  width=0.88\textwidth, height=0.68\textheight,
  xlabel={Fiscal Year (ending March)},
  ylabel={Federal Budget Balance (\$billions)},
  xmin=1990, xmax=2025,
  ymin=-340, ymax=20,
  xtick={1990,1995,2000,2005,2010,2015,2020,2025},
  xticklabel style={rotate=45, anchor=east, font=\tiny},
  grid=both, grid style={dotted,gray!40},
  tick label style={font=\tiny},
  label style={font=\small},
  title style={font=\small\bfseries},
  title={Canada's Federal Budget Balance: Deficits Dominate (Surplus Era 1998--2008)}
]
\addplot[canadared, very thick, ybar, bar width=0.35] coordinates {
  (1990,-12)(1991,-30)(1992,-35)(1993,-40)(1994,-36)(1995,-30)
  (1996,-20)(1997,-8)(1998,3)(1999,7)(2000,12)(2001,15)
  (2002,8)(2003,7)(2004,9)(2005,13)(2006,14)(2007,10)(2008,1)
  (2009,-55)(2010,-33)(2011,-25)(2012,-18)(2013,-12)(2014,-5)
  (2015,-2)(2016,-6)(2017,-15)(2018,-19)(2019,-14)(2020,-314)
  (2021,-90)(2022,-10)(2023,-36)(2024,-40)
};
\addplot[goldenrod, dashed, thick] coordinates {(1990,0)(2025,0)};
\node[font=\tiny, fill=leafgreen!10, draw=leafgreen, rounded corners, align=center]
  at (axis cs:2003,15) {Surplus years\\1998--2008};
\node[font=\tiny, fill=lightred, draw=canadared, rounded corners, align=center]
  at (axis cs:2020,-280) {COVID-19:\\$-$\$314B};
\end{axis}
\end{tikzpicture}
\end{center}
{\tiny Source: Department of Finance Canada, Federal Budget. Positive = surplus; negative = deficit.}
\end{frame}

\begin{frame}{Automatic Stabilizers and the Business Cycle}
\transfade
\begin{itemize}[<+->]
  \item \textbf{In recessions} (automatic stabilizers kick in):
  \begin{itemize}[<+->]
    \item Tax revenues \alert{fall} (lower incomes, profits, consumption).
    \item EI payments, social assistance \alert{rise}.
    \item Result: budget deficit grows automatically, helping support AD.
  \end{itemize}
  \item \textbf{In expansions}:
  \begin{itemize}[<+->]
    \item Tax revenues \alert{rise}.
    \item EI, welfare payments \alert{fall}.
    \item Result: budget moves toward surplus automatically, dampening inflation.
  \end{itemize}
  \item \textbf{Canada's EI example (2020):}
  \begin{itemize}[<+->]
    \item EI and emergency benefits paid out $\approx \$100$ billion in 2020 alone.
    \item This replaced lost income for millions of Canadians, supporting consumption.
    \item Without this, the recession would have been \textbf{much deeper}.
  \end{itemize}
\end{itemize}
\end{frame}

\begin{frame}{Figure 12.6 -- Federal Government Debt to GDP Ratio (Canada)}
\begin{center}
\begin{tikzpicture}
\begin{axis}[
  width=0.88\textwidth, height=0.68\textheight,
  xlabel={Year},
  ylabel={Federal Debt-to-GDP Ratio (\%)},
  xmin=1980, xmax=2025,
  ymin=0, ymax=75,
  xtick={1980,1985,1990,1995,2000,2005,2010,2015,2020,2025},
  xticklabel style={rotate=45, anchor=east, font=\tiny},
  grid=both, grid style={dotted,gray!40},
  tick label style={font=\tiny},
  label style={font=\small},
  title style={font=\small\bfseries},
  title={Federal Debt-to-GDP Rose Sharply in 1990s, Then Fell, Then Rose Again (2020)}
]
\addplot[accent, very thick, mark=none] coordinates {
  (1980,25)(1985,40)(1990,50)(1993,68)(1996,68)(2000,55)
  (2005,35)(2008,28)(2010,33)(2015,30)(2018,27)(2019,27)
  (2020,47)(2021,50)(2022,43)(2023,42)(2024,42)
};
\addplot[goldenrod, dashed, thick] coordinates {(1980,45)(2025,45)}
  node[right, font=\tiny, color=goldenrod] {G7 avg.};
\node[font=\tiny, fill=lightred, draw=canadared, rounded corners, align=center]
  at (axis cs:1995,64) {1993--1996 peak:\\Moody's threatened\\to downgrade Canada};
\node[font=\tiny, fill=leafgreen!10, draw=leafgreen, rounded corners, align=center]
  at (axis cs:2010,20) {Sustained decline\\2000--2019};
\node[font=\tiny, fill=lightred, draw=canadared, rounded corners, align=center]
  at (axis cs:2021,58) {COVID-19\\spike};
\end{axis}
\end{tikzpicture}
\end{center}
{\tiny Source: Department of Finance Canada; IMF Fiscal Monitor. G7 average is approximate.}
\end{frame}

\begin{frame}{Is Government Debt a Problem?}
\transfade
\begin{itemize}[<+->]
  \item \textbf{Short-term view:} Canada's federal debt-to-GDP ratio ($\approx$42\%) is lower than most G7 countries. Not an immediate crisis.
  \item \textbf{Long-term concerns:}
  \begin{itemize}[<+->]
    \item If debt grows faster than GDP, interest payments consume a larger share of the budget.
    \item \textbf{Crowding out:} rising debt may gradually push up real interest rates, reducing private investment.
    \item \textbf{Generational equity:} future taxpayers (today's children) inherit the debt.
  \end{itemize}
  \item \textbf{When is debt less problematic?}
  \begin{itemize}[<+->]
    \item When it finances \textbf{productive investments}: infrastructure, education, R\&D.
    \item These raise future potential GDP, making the debt easier to service.
    \item \textbf{Canada 2020:} emergency spending was necessary to prevent economic collapse, even though it was largely consumption-focused.
  \end{itemize}
\end{itemize}
\end{frame}

% ============================================================
\section{12.7 Fiscal Policy in the Long Run}
% ============================================================

\begin{frame}{12.7 Learning Objective}
\transfade
\begin{block}{Goal}
Discuss the effects of fiscal policy in the long run (supply-side effects).
\end{block}
\begin{itemize}[<+->]
  \item So far: fiscal policy as a short-run stabilization tool (shifts AD).
  \item Now: how fiscal policy affects \textbf{potential GDP} (shifts LRAS).
  \item \textbf{Supply-side economics:}
  \begin{itemize}[<+->]
    \item Fiscal policy actions aimed at increasing long-run aggregate supply.
    \item Focus: taxes and regulations that affect incentives to work, save, invest, and innovate.
  \end{itemize}
\end{itemize}
\end{frame}

\begin{frame}{The Tax Wedge and Incentives}
\transfade
\begin{itemize}[<+->]
  \item Most taxes create a \textbf{wedge} between the pre-tax and post-tax return to an activity.
  \item \textbf{Examples:}
  \begin{itemize}[<+->]
    \item \textbf{Personal income tax:} if your marginal rate is 53\% (federal + BC combined), each extra dollar earned yields only 47 cents after tax.
    \begin{itemize}[<+->]
      \item May reduce hours worked, entrepreneurship, or human capital investment.
    \end{itemize}
    \item \textbf{Corporate income tax:} Canada's combined rate $\approx$26\% in 2024.
    \begin{itemize}[<+->]
      \item High rates may reduce investment in new machinery, R\&D, and job creation.
    \end{itemize}
    \item \textbf{Capital gains tax:} affects saving and investment decisions.
    \begin{itemize}[<+->]
      \item \textbf{2024 Canada:} Trudeau government raised capital gains inclusion rate from 50\% to 67\% --- controversial impact on investment incentives.
    \end{itemize}
  \end{itemize}
\end{itemize}
\end{frame}

\begin{frame}{Figure 12.7 -- Supply-Side Effects of a Tax Cut}
\begin{center}
\begin{tikzpicture}
\begin{axis}[
  width=0.80\textwidth, height=0.70\textheight,
  xlabel={Real GDP (\$billions)},
  ylabel={Price Level ($P$)},
  xmin=0, xmax=3500, ymin=80, ymax=200,
  xtick={500,1000,1500,2000,2500,3000,3500},
  ytick={90,100,110,120,130,140,150,160,170,180,190},
  grid=major, grid style={dotted,gray!40},
  tick label style={font=\small},
  label style={font=\small},
  title style={font=\small\bfseries},
  clip=false,
  title={Supply-Side Tax Reform: LRAS Shifts Right, Price Level Falls}
]
% Consistent equations:
% AD: P=170-0.02x
% SRAS1: P=0.03x+70  (at x=2000: P=130) -- old LR eq at (2000,130) with AD
% LRAS1 at x=2000  [AD=SRAS1=130 at x=2000 -- consistent triple intersection]
% After reform:
% SRAS2: P=0.03x+45  (shifted down -- lower costs)
% LRAS2 at x=2500  [AD=SRAS2=120 at x=2500 -- consistent triple intersection]
% Verification: AD at 2500: 170-0.02*2500=120. SRAS2 at 2500: 0.03*2500+45=120. ✓

% LRAS1 (dashed)
\addplot[shiftpurple!50, very thick, dashed] coordinates {(2000,80)(2000,200)};
\node[above, font=\tiny, color=shiftpurple!60] at (axis cs:2000,203) {$LRAS_1$};
% LRAS2 (solid, right of LRAS1)
\addplot[shiftpurple, very thick] coordinates {(2500,80)(2500,200)};
\node[above, font=\tiny, color=shiftpurple] at (axis cs:2500,203) {$LRAS_2$};
% SRAS1 (dashed)
\addplot[goldenrod!60, very thick, dashed, domain=300:3200] {0.03*x + 70};
\node[right, font=\tiny, color=goldenrod!60] at (axis cs:3050,161) {$SRAS_1$};
% SRAS2 (solid, shifted right/down)
\addplot[goldenrod, very thick, domain=300:3400] {0.03*x + 45};
\node[right, font=\tiny, color=goldenrod] at (axis cs:3200,141) {$SRAS_2$};
% AD: P=170-0.02x
\addplot[deeporange, very thick, domain=300:3400] {170 - 0.02*x}
  node[right, font=\tiny, color=deeporange] {\textbf{AD}};
% Old LR equilibrium: (2000, 130)
\addplot[accent, only marks, mark=*, mark size=4pt] coordinates {(2000,130)};
\node[above left, font=\tiny, color=accent] at (axis cs:2000,130) {$E_1$};
\draw[dashed, accent!60] (axis cs:2000,130)--(axis cs:0,130)
  node[left, font=\tiny, color=accent] {$P_1=130$};
% New LR equilibrium: (2500, 120)
\addplot[leafgreen, only marks, mark=*, mark size=4pt] coordinates {(2500,120)};
\node[above right, font=\tiny, color=leafgreen] at (axis cs:2500,120) {$E_2$};
\draw[dashed, leafgreen!60] (axis cs:2500,120)--(axis cs:0,120)
  node[left, font=\tiny, color=leafgreen] {$P_2=120$};
% Arrow showing expansion of potential GDP
\draw[-{Stealth[length=7pt]}, very thick, shiftpurple]
  (axis cs:2050,170) -- (axis cs:2450,170)
  node[midway, above, font=\tiny, color=shiftpurple] {Potential GDP rises};
\node[font=\tiny, fill=leafgreen!10, draw=leafgreen, rounded corners, align=center]
  at (axis cs:850,185) {Tax reform $\Rightarrow$ more work \& investment\\$\Rightarrow$ LRAS and SRAS shift right\\$\Rightarrow$ higher $Y^*$, lower $P$};
\end{axis}
\end{tikzpicture}
\end{center}
{\footnotesize Supply-side tax reform raises potential GDP from \$2{,}000B to \$2{,}500B while lowering the long-run price level.}
\end{frame}

\begin{frame}{Long-Run Effects of Tax Reform: What We Know}
\transfade
\begin{itemize}[<+->]
  \item Tax reform can increase potential GDP in the long run by:
  \begin{itemize}[<+->]
    \item Raising labour supply (more hours worked, more participation).
    \item Increasing saving and investment (more capital formation).
    \item Encouraging entrepreneurship and innovation.
  \end{itemize}
  \item \textbf{Uncertainty:}
  \begin{itemize}[<+->]
    \item Magnitude of supply-side effects is \textbf{empirically contested}.
    \item People may want to work more at lower marginal rates, but may face \textit{constraints} (fixed schedules, childcare).
    \item Tax cuts also create short-run demand effects (shifting AD right), which may dominate.
  \end{itemize}
  \item \textbf{Canadian policy debate:}
  \begin{itemize}[<+->]
    \item 2024 capital gains inclusion rate increase: critics argued it would deter investment; supporters argued revenue would fund social priorities.
    \item Active empirical research ongoing.
  \end{itemize}
\end{itemize}
\end{frame}

% ============================================================
\section*{Appendix: A Closer Look at the Multiplier}
% ============================================================

\begin{frame}{Appendix -- The Government Purchases and Tax Multipliers}
\transfade
\begin{itemize}[<+->]
  \item Simple Canadian macroeconomic model (closed economy, fixed price level):
  \begin{align*}
    C &= 1{,}000 + 0.8(Y - T), \quad T = 250, \\
    I &= 300, \quad G = 500.
  \end{align*}
  \item \textbf{Equilibrium:} $Y = C + I + G$:
  \[
    Y = 1{,}000 + 0.8(Y - 250) + 300 + 500.
  \]
  \[
    Y = 1{,}600 + 0.8Y - 200 = 1{,}400 + 0.8Y.
  \]
  \[
    0.2Y = 1{,}400 \quad \Rightarrow \quad Y = \$7{,}000\text{ billion}.
  \]
  \item \textbf{Government purchases multiplier:} $\frac{\Delta Y}{\Delta G} = \frac{1}{1 - 0.8} = 5$.
  \item $\Delta G = +\$10B \Rightarrow \Delta Y = +\$50B$.
\end{itemize}
\end{frame}

\begin{frame}{The Tax Multiplier and the Balanced Budget Multiplier}
\transfade
\begin{itemize}[<+->]
  \item \textbf{Tax multiplier:}
  \[
    \frac{\Delta Y}{\Delta T} = -\frac{\text{MPC}}{1 - \text{MPC}} = -\frac{0.8}{0.2} = -4.
  \]
  \item $\Delta T = -\$10B$ (tax cut) $\Rightarrow \Delta Y = +\$40B$.
  \item \textbf{Why smaller than $G$ multiplier?} A \$10B tax cut gives \$10B to households, but they \textbf{save} $0.2 \times 10 = \$2B$, so only \$8B is initially spent.
  \item \textbf{Balanced budget multiplier:}
  \begin{itemize}[<+->]
    \item Increase $G$ and $T$ by the same amount (\$10B).
    \item $\Delta Y = (+5) \times 10B + (-4) \times 10B = \$10B$.
    \item Balanced budget multiplier $= \frac{\Delta Y}{\Delta G} = 1$.
    \item Equal dollar increases in spending and taxes raise GDP by the same dollar.
  \end{itemize}
\end{itemize}
\end{frame}

% ============================================================
\section*{Chapter Summary}
% ============================================================

\begin{frame}{Summary}
\transfade
\begin{itemize}[<+->]
  \item \textbf{Fiscal policy:} changes in federal spending ($G$) and taxes ($T$) to achieve macroeconomic goals.
  \item \textbf{Automatic stabilizers} (EI, progressive income tax) cushion downturns without legislation.
  \item \textbf{Expansionary} ($\uparrow G$ or $\downarrow T$) closes recessionary gaps; \textbf{contractionary} ($\downarrow G$ or $\uparrow T$) reduces inflationary pressure.
  \item Government purchases multiplier $= \frac{1}{1-\text{MPC}}$; tax multiplier $= -\frac{\text{MPC}}{1-\text{MPC}}$.
  \item Practical limits: \textbf{timing lags} (recognition, legislative, implementation) and \textbf{crowding out}.
  \item Canada's federal debt-to-GDP $\approx$42\% (2024) --- manageable but trending higher.
  \item \textbf{Supply-side} fiscal policy (tax reform) can shift LRAS right and raise potential GDP, but magnitudes are uncertain.
  \item \textbf{Balanced budget multiplier $= 1$}: equal increases in $G$ and $T$ raise GDP by the same dollar.
  \item \textbf{Canada 2020 case study:} \$314B in emergency spending was the largest ever fiscal expansion --- essential to prevent depression-level output losses, at the cost of a large deficit and higher debt.
\end{itemize}
\end{frame}

\begin{frame}
\centering
\Huge \textbf{Thank you!}\\[1em]
\large Questions?
\end{frame}

\end{document}
